% Source: https://artofproblemsolving.com/community/c3h1902252_counting
% Stable identity: aops-topic-1902252-q01
Sort the lines by direction: lines in the same class are parallel, and lines in different classes meet in exactly one point. Let the class sizes be a partition of $4$.
\begin{itemize}
\item $\{4\}$: all parallel, no intersection point at all.
\item $\{2,2\}$: each line of one pair meets each line of the other pair, in $4$ different points (the vertices of a parallelogram). Too many.
\item $\{1,1,1,1\}$: no two lines are parallel, so there are $6$ intersecting pairs. If a point lies on $m$ lines, it accounts for $\dbinom m2\in\{1,3,6\}$ of these pairs. Three points would need three numbers from $\{1,3,6\}$ adding to $6$, which is impossible. So this case cannot have exactly three points.
\item $\{3,1\}$: three parallel lines and one line crossing them in $3$ different points. This works.
\item $\{2,1,1\}$: a parallel pair $\ell_1,\ell_2$ and two more lines $m_1,m_2$ meeting at a point $Q$. The points $\ell_i\cap m_j$ are four different points unless $Q$ lies on $\ell_1$ or $\ell_2$. If $Q$ is on neither, there are $5$ points. If $Q$ lies on $\ell_1$ (say), the points are $Q$, $\ell_2\cap m_1$ and $\ell_2\cap m_2$, exactly three. This works.
\end{itemize}

\textbf{Counting regions.} A new line that meets the earlier lines in $c$ distinct points is cut into $c+1$ pieces, and each piece splits a region in two.

\emph{Case $\{3,1\}$:} three parallel lines make $4$ strips, and the crossing line adds $3+1=4$ regions: $8$ regions.

\begin{asy}
usepackage("amsmath");
size(5cm);
defaultpen(fontsize(10pt));
draw((-3,0)--(3,0),Arrows(TeXHead));draw((-3,1.5)--(3,1.5),Arrows(TeXHead));draw((-3,3)--(3,3),Arrows(TeXHead));
draw((-2,-1)--(2.4,4.5),blue,Arrows(TeXHead));
pen q=fontsize(9pt);
label("$1$",(-2,3.8),q);label("$2$",(2.3,3.8),q);label("$3$",(-1.6,2.25),q);label("$4$",(1.9,2.25),q);
label("$5$",(-2.2,0.75),q);label("$6$",(1.5,0.75),q);label("$7$",(-2.6,-0.7),q);label("$8$",(0.6,-0.7),q);
label("\mbox{}",(min(currentpicture,user=true)+max(currentpicture,user=true))/2,black); // keeps the text after the figure black when LaTeX typesets the labels inline
\end{asy}

\emph{Case $\{2,1,1\}$ with $Q\in\ell_1$:} the three lines through $Q$ make $6$ regions, and $\ell_2$ crosses the other two lines at $2$ points, adding $3$ regions: $9$ regions.

\begin{asy}
usepackage("amsmath");
size(5.5cm);
defaultpen(fontsize(10pt));
draw((-3.5,0)--(3.5,0),Arrows(TeXHead));draw((-3.5,2)--(3.5,2),Arrows(TeXHead));
draw((-1.8,-2.16)--(2.9,3.48),blue,Arrows(TeXHead));draw((1.8,-2.16)--(-2.9,3.48),blue,Arrows(TeXHead));
dot((0,0));label("$Q$",(0,0),3*S);
pen q=fontsize(9pt);
label("$1$",(0,3.0),q);label("$2$",(-3.0,2.8),q);label("$3$",(3.0,2.8),q);label("$4$",(0,1.2),q);
label("$5$",(-2.6,1.0),q);label("$6$",(2.6,1.0),q);label("$7$",(0,-1.5),q);label("$8$",(-2.6,-1.0),q);label("$9$",(2.6,-1.0),q);
label("\mbox{}",(min(currentpicture,user=true)+max(currentpicture,user=true))/2,black); // keeps the text after the figure black when LaTeX typesets the labels inline
\end{asy}

So the number of regions is $\boxed{8\text{ or }9}$, and both values occur. (The answer $13$ is impossible: four lines can never make more than $1+1+2+3+4=11$ regions.)
